\documentclass[11pt,a4paper]{article}
\usepackage[margin=2.5cm]{geometry}
\usepackage{amsmath,amssymb}
\usepackage{booktabs}
\usepackage{parskip}
\usepackage{hyperref}
\usepackage{xcolor}
\usepackage{graphicx}
\usepackage{siunitx}

\hypersetup{colorlinks=true,linkcolor=blue,citecolor=blue,urlcolor=blue}

\newcommand{\DS}{\Delta\Sigma}
\newcommand{\Msun}{\,M_{\odot}}
\newcommand{\Mpc}{\,\mathrm{Mpc}}
\newcommand{\gtrue}{\gamma_{\rm true}}

\title{\textbf{Spherical vs.\ Triaxial NFW $\DS$ Profiles}\\[0.4em]
\large Controlled comparison at zero shape noise --- destacking project}
\author{destacking project}
\date{2026 May 26}

\begin{document}
\maketitle

\begin{abstract}
We measure the lensing $\DS$ signal of identical NFW haloes
($z=0.4$, $M_{200}=10^{14}\Msun$, $c=5$) under two different shear-application
schemes from \texttt{xlensing.sims}: spherical
(\texttt{apply\_NFW\_shear\_region}) and triaxial
(\texttt{apply\_triaxial\_NFW\_shear\_region}) with shapes drawn from the
prior used in \texttt{gamma\_scan.py} ($b/a\sim U(0.5,1)$,
$c/a\sim U(0.3,b/a)$, isotropic line of sight).  Each cluster uses the
\emph{same} zero-shape-noise background-galaxy field so the only difference
in the recovered $\DS$ is the geometric projection of the ellipsoid.
The orientation-averaged triaxial signal is $14$--$17\,\%$ \emph{lower} than
the spherical signal at every radius in our fit range, with strong
LOS-orientation dependence: $-35\,\%$ for edge-on, $-11\,\%$ for face-on.
This is exactly the systematic that drives the
$\hat\gamma\approx-0.5$ plateau in the gamma scan.
\end{abstract}

\tableofcontents
\newpage

% ============================================================
\section{Setup}
% ============================================================

We fix a single cluster identity and vary only the triaxial shape.  The
spherical and triaxial measurements differ \emph{only} in the shear-application
function; everything else (mass, concentration, redshift, background galaxies,
binning) is identical.

\begin{table}[htbp]
  \centering
  \begin{tabular}{ll}
    \toprule
    parameter & value\\
    \midrule
    cluster redshift $z$               & $0.4$\\
    cluster mass $M_{200}$             & $1.0\times 10^{14}\Msun$\\
    cluster concentration $c$          & $5.0$\\
    number of shape draws $N$          & $100$\\
    background galaxy density          & $30$\,gal/arcmin$^2$\\
    field half-width                   & $0.006$\,rad ($\approx 20.6$\,arcmin)\\
    source redshift range              & $U(0.6,\,2.5)$\\
    shape noise                        & $0.0$ (intrinsic $e_1=e_2=0$)\\
    metacal response $R_{ii}$          & $1$ (no calibration correction)\\
    radial bins                        & $12$, log-spaced in $[0.16,\,6.3]\Mpc$\\
    \bottomrule
  \end{tabular}
  \caption{Controlled-experiment configuration.  All clusters share
  $(z,M,c)$; only the triaxial shape and LOS direction vary.}
  \label{tab:config}
\end{table}

For each cluster we:
\begin{enumerate}
\item sample shape parameters $(b/a,\,c/a,\,\theta_{\rm LOS},\,\phi_{\rm LOS})$
      from the gamma-scan prior;
\item generate one zero-shape-noise background-galaxy catalogue;
\item apply spherical NFW shear to a copy of that catalogue and measure
      $\DS_{\rm sph}(R)$ via the same
      \texttt{metacal\_cluster\_lensing}\,$\to$\,\texttt{signal} pipeline used
      by \texttt{gamma\_scan.py};
\item apply triaxial NFW shear to a second copy of the same catalogue and
      measure $\DS_{\rm tri}(R)$.
\end{enumerate}

The script is \texttt{scripts/sphere\_vs\_triaxial.py}; numerical outputs are
stored in \texttt{sphere\_vs\_triaxial\_outputs/}.  Because there is no shape
noise and the response matrix is unity, the recovered $\DS$ in every bin is
the (numerically-binned) projection of the NFW profile times the critical
surface density of the source plane --- no random scatter, no calibration
fudge factor.

% ============================================================
\section{Population-mean profiles}
% ============================================================

Figure~\ref{fig:profiles} overlays each cluster's triaxial $\DS$ profile (pale
blue lines, $N=100$) and shows the population means for both schemes.  The
spherical population mean (black) follows the canonical NFW shape; the
triaxial population mean (steelblue) is consistently below it across every
radial bin.

\begin{figure}[htbp]
  \centering
  \includegraphics[width=0.78\linewidth]{../sphere_vs_triaxial_outputs/profiles.png}
  \caption{$\DS(R)$ at $z=0.4$, $M=10^{14}\Msun$, $c=5$ for spherical NFW
  (black) and triaxial NFW (steelblue), averaged over $N=100$ shape draws.
  Pale lines: individual triaxial realisations.  The triaxial population
  signal is uniformly lower than the spherical reference.}
  \label{fig:profiles}
\end{figure}

Table~\ref{tab:perR} gives the numbers per radial bin: the per-cluster
\emph{median} fractional residual is $-14$ to $-17\,\%$, the \emph{mean}
fractional residual is $-18$ to $-20\,\%$ (the mean is more negative than the
median because the per-cluster distribution has a long negative tail from
edge-on prolate orientations), and the 68\,\% per-cluster band ranges from
$\sim-34\,\%$ to $\sim-3\,\%$.  The systematic is clearly resolved at the
single-bin level.

\begin{table}[htbp]
  \centering
  \small
  \begin{tabular}{rrrrr}
    \toprule
    $R$ [Mpc] & $\langle\DS_{\rm sph}\rangle$ & $\langle\DS_{\rm tri}\rangle$
              & median frac.\ resid.\ & 68\,\% band\\
    \midrule
    $0.158$ & $91.82$ & $74.24$ & $-14.0\,\%$ & $[-33.3,\,-2.7]\,\%$\\
    $0.222$ & $78.85$ & $64.41$ & $-14.7\,\%$ & $[-33.9,\,-4.0]\,\%$\\
    $0.310$ & $64.33$ & $52.13$ & $-14.6\,\%$ & $[-30.5,\,-4.3]\,\%$\\
    $0.433$ & $50.22$ & $40.03$ & $-17.0\,\%$ & $[-33.2,\,-3.3]\,\%$\\
    $0.605$ & $37.47$ & $30.08$ & $-15.8\,\%$ & $[-34.3,\,-4.0]\,\%$\\
    $0.846$ & $26.72$ & $21.36$ & $-17.3\,\%$ & $[-33.5,\,-3.9]\,\%$\\
    $1.182$ & $18.31$ & $14.60$ & $-17.3\,\%$ & $[-34.8,\,-3.6]\,\%$\\
    $1.653$ & $12.07$ & $\phantom{0}9.65$ & $-16.5\,\%$ & $[-35.4,\,-3.5]\,\%$\\
    $2.310$ & $\phantom{0}7.72$ & $\phantom{0}6.19$ & $-16.7\,\%$ & $[-34.6,\,-3.8]\,\%$\\
    $3.229$ & $\phantom{0}4.80$ & $\phantom{0}3.84$ & $-16.7\,\%$ & $[-34.0,\,-3.5]\,\%$\\
    $4.514$ & $\phantom{0}2.92$ & $\phantom{0}2.34$ & $-16.7\,\%$ & $[-34.6,\,-3.5]\,\%$\\
    $6.310$ & $\phantom{0}1.87$ & $\phantom{0}1.50$ & $-17.0\,\%$ & $[-34.4,\,-3.4]\,\%$\\
    \bottomrule
  \end{tabular}
  \caption{Per-radius mean $\DS$ (in $M_\odot/\mathrm{pc}^2$) under spherical
  and triaxial shear, with the per-cluster fractional residual statistics.
  The deficit is essentially radius-independent.}
  \label{tab:perR}
\end{table}

% ============================================================
\section{Fractional residual}
% ============================================================

Figure~\ref{fig:frac} plots the per-cluster fractional residual
$(\DS_{\rm tri}-\DS_{\rm sph})/\DS_{\rm sph}$ as a function of $R$.  The
median (steelblue line) is flat at $-15$ to $-17\,\%$, the mean (black
dashed) is similarly flat at $-19$ to $-20\,\%$, and the 68\,\% per-cluster
band (shaded) spans roughly $-35\,\%$ to $-3\,\%$.  The flatness across all
radii is the central empirical observation: the triaxial vs.\ spherical
mismatch is \emph{not} a localised inner- or outer-profile effect; it is
a near-uniform amplitude shift.

\begin{figure}[htbp]
  \centering
  \includegraphics[width=0.78\linewidth]{../sphere_vs_triaxial_outputs/fractional_residual.png}
  \caption{Fractional residual
  $(\DS_{\rm tri}-\DS_{\rm sph})/\DS_{\rm sph}$ vs.\ $R$.  Shaded:
  $16$--$84\,\%$ per-cluster band; solid steelblue: median; dashed black:
  mean.  The systematic is flat in radius and asymmetric in distribution
  (mean below median).}
  \label{fig:frac}
\end{figure}

% ============================================================
\section{Orientation dependence}
% ============================================================

The dominant driver of the per-cluster variance is the line-of-sight (LOS)
orientation.  Splitting the population by $\cos\theta_{\rm LOS}$ (with
$\theta=0$ being the major axis along the LOS, ``face-on'')
gives Fig.~\ref{fig:byorient}:

\begin{itemize}
  \item \textbf{Face-on}, $\cos\theta\ge 0.66$ ($N=33$):
        fractional residual $\approx -11\,\%$ at all radii.  The major axis
        is roughly aligned with the LOS; the projected ellipsoid is
        compact and the projected $\DS$ is closest to a spherical signal of
        the same intrinsic concentration.
  \item \textbf{Intermediate}, $0.33\le\cos\theta < 0.66$ ($N=38$):
        $\approx -15\,\%$.
  \item \textbf{Edge-on}, $\cos\theta<0.33$ ($N=29$):
        $\approx -35\,\%$.  The major axis lies near the plane of the sky;
        most of the mass projects to large radii in the azimuthal average,
        and the apparent inner profile drops sharply relative to spherical.
\end{itemize}

\begin{figure}[htbp]
  \centering
  \includegraphics[width=0.78\linewidth]{../sphere_vs_triaxial_outputs/by_orientation.png}
  \caption{Population-mean fractional residual split by LOS orientation.
  Edge-on triaxial halos produce $\DS$ profiles $\sim 35\,\%$ below the
  spherical reference; face-on halos differ by $\sim 11\,\%$.  Each bin
  is flat in radius.}
  \label{fig:byorient}
\end{figure}

The factor of $\sim 3$ between face-on and edge-on is the largest single
source of inter-cluster variance in the deficit.  Because $\cos\theta_{\rm
LOS}$ is uniformly distributed on $[0,1]$ for isotropic LOS, all three bins
have comparable populations ($\sim$~33 each), so the population mean lands
firmly inside the intermediate regime.

% ============================================================
\section{Dependence on intrinsic triaxiality $T$}
% ============================================================

We classify each halo by the Franx+91 triaxiality parameter
$T=(a^2-b^2)/(a^2-c^2)$:
oblate ($T<1/3$), generally triaxial ($1/3\le T\le 2/3$), prolate
($T>2/3$).  Splitting by $T$ at fixed mass and concentration gives
Fig.~\ref{fig:bytri}:

\begin{itemize}
  \item oblate ($N=21$): residual $\approx -20\,\%$
  \item triaxial ($N=34$): residual $\approx -19\,\%$
  \item prolate ($N=45$): residual $\approx -21\,\%$
\end{itemize}

\begin{figure}[htbp]
  \centering
  \includegraphics[width=0.78\linewidth]{../sphere_vs_triaxial_outputs/by_triaxiality.png}
  \caption{Population-mean fractional residual split by intrinsic
  triaxiality $T$.  The three categories agree to within $\sim 2\,\%$ ---
  much smaller than the orientation effect.}
  \label{fig:bytri}
\end{figure}

The three categories are statistically indistinguishable.  This is the
expected pattern: \emph{intrinsic} triaxiality sets the geometry of the
ellipsoid, but the lensing signal depends on the \emph{projected} ellipse
on the sky, which is set jointly by intrinsic shape and orientation.  With
isotropic LOS averaging, the orientation absorbs most of the
shape-information and the intrinsic-$T$ split flattens out.

% ============================================================
\section{Connection to the $\hat\gamma$ plateau}
% ============================================================

The diagnosis of the gamma plateau in
\texttt{2026-05-26\_gamma\_plateau\_diagnosis.pdf} argued that the
$\hat\gamma\approx -0.5$ floor for negative $\gamma_{\rm true}$ comes from
the spherical-vs.-triaxial mismatch, which is invariant under $c(M,z,\Delta
m)$ reparametrisation.  The numbers in this report sharpen that statement.

\begin{enumerate}
\item At fixed mass and concentration, the orientation-averaged triaxial
      $\DS$ is $\sim 18$--$20\,\%$ \emph{lower} than the spherical NFW
      reference (mean over 100 clusters at $c=5$, $M=10^{14}\Msun$).
\item The deficit is \emph{nearly flat in radius}, so the spherical model
      cannot recover it by warping the radial shape; it can only do so by
      changing the overall normalisation, which in this pipeline means
      changing $(c_0,\alpha,\beta,\gamma,\log M_{\rm norm})$.
\item Of those five parameters, $c_0$, $\alpha$, $\beta$, and
      $\log M_{\rm norm}$ are tightly priored at their Dutton--Maccio
      / mass-richness calibration values.  Only $\gamma$ has prior width
      $3.5$.  So $\gamma$ is forced to absorb the
      triaxial-vs.-spherical amplitude deficit population-wide.
\item Negative $\gtrue$ pushes the implied $c$-distribution bimodal: some
      clusters at $c\!\gg\!20$, others at $c\!\ll\!1$.  At $c=5$ the
      triaxial-vs.-spherical mismatch already costs $\sim 18\,\%$; at
      extreme $c$ the spherical projection becomes a worse approximation
      to the orientation-averaged triaxial signal.  The likelihood
      therefore favours the $c$-distribution that stays bounded in
      $\sim[1,8]$, which in the current parametrisation corresponds to
      $\gamma\approx -0.5$, exactly the plateau location.
\end{enumerate}

The $\sim 20\,\%$ amplitude offset measured here is not a physically
``small'' number for a $10^4$-cluster joint fit at percent-level shape-noise
sensitivity --- it is the dominant systematic.  The fact that it lives
entirely in the geometry, and is invariant under the parametric
$c$-relation, is why no $c$-relation reparametrisation can move the
plateau.

% ============================================================
\section{Caveats}
% ============================================================

\textbf{Normalisation convention.}\quad
The $\sim 15$--$20\,\%$ uniform-in-radius deficit
is partly tied to how \texttt{xlensing.sims.apply\_triaxial\_NFW\_shear\_region}
defines the triaxial halo's enclosed mass.  The function takes
$(M_{200},\,c,\,a,b,c_{\rm axis})$ with $a=r_{\rm vir}(M_{200})$ but $b,c<a$:
the iso-density ellipsoid is then smaller in volume than the equivalent
sphere of radius $r_{\rm vir}$, so the orientation-averaged projected
amplitude can naturally come out lower at fixed $(M_{200},c)$.  This is the
relevant convention for the destacking pipeline because it is the one used
\emph{both} in the simulation and in the comparison here.  But it means the
absolute level of the deficit is not a universal physical statement --- it
is the deficit \emph{as defined by the pipeline's own simulator}.

\textbf{Fixed cluster identity.}\quad
We hold $(z,M,c)$ fixed.  In real cluster samples there will be a
mass--orientation correlation through selection, and concentration scatter
on top of the shape scatter; both would broaden the per-cluster
distribution.  This test isolates the geometric piece.

\textbf{No miscentering, no concentration scatter.}\quad
The gamma\_scan simulations also include cluster miscentering (\,$F_{\rm
cen}=0.64$, Rayleigh $\sigma_{\rm off}=0.32\,\mathrm{Mpc}/h$) and concentration
scatter inherited from the \texttt{generate\_truth\_catalog} pipeline.  Both
are turned off here so that the only effect we see is triaxiality.
Miscentering would suppress the inner $\sim 1\Mpc$ further; concentration
scatter would widen the per-cluster band.

\textbf{Real data carry no triaxiality information.}\quad
This is the headline practical implication.  Even with a perfect
characterisation of the deficit shown here, the destacking pipeline cannot
absorb it through any per-cluster nuisance because there is no per-cluster
observable that recovers shape or orientation.  Population-level priors are
possible in principle but require external calibration (N-body simulations,
projected-axis-ratio statistics from richness/X-ray morphology) that the
real survey data the pipeline is being designed for do not supply.

% ============================================================
\section*{Files}
% ============================================================

\begin{itemize}
  \item \texttt{scripts/sphere\_vs\_triaxial.py} --- the comparison driver.
  \item \texttt{sphere\_vs\_triaxial\_outputs/profiles.npz} --- per-cluster
        $\DS_{\rm sph}$, $\DS_{\rm tri}$ and shape parameters.
  \item \texttt{sphere\_vs\_triaxial\_outputs/summary.json} --- numerical
        summary (per-bin means, medians, std, fractional-residual
        percentiles).
  \item Plots referenced in this report:
        \texttt{profiles.png},
        \texttt{fractional\_residual.png},
        \texttt{by\_orientation.png},
        \texttt{by\_triaxiality.png}.
\end{itemize}

\end{document}
